\section{Conclusion}
\label{sec:conclusion}

This paper answers the adversarial reframing of B.17's parameter-sensitivity
finding: the maximum partisan shift achievable by optimizing all five
\bisect{} parameters for partisan advantage is 0.3 Democratic seats nationally.

This is a tight bound for three structural reasons:
(1)~parameters affect leaf-level cut quality within a fixed bisection tree,
not the tree structure itself;
(2)~competitive-district geography is not sensitive to 1--3 tract boundary
shifts in non-competitive regions;
and (3)~the five parameters interact on the same geometric bottleneck,
producing saturation rather than amplification.

The DIA pre-registration protocol eliminates this gaming vector entirely
by requiring all parameters to be fixed and publicly recorded before Census
redistricting data are released.
Under pre-registration, the question ``could they have set parameters to
favor their party?'' is converted from an unprovable allegation into a
verifiable cryptographic fact: either the deployed parameters match the
pre-registered hash or they do not.

The broader lesson for algorithmic redistricting design is that parameter
gaming is a minor concern relative to algorithm-structure gaming (B.0's
18-seat gap) and geographic-definition gaming (R.3's 1.5-seat gap).
The DIA should concentrate its structural constraints on algorithm
pre-specification and geographic definition mandates, treating parameter
pre-registration as a lower-stakes but still necessary safeguard.
